\documentclass[]{article}
\usepackage{graphicx} % Required for inserting images
\usepackage{amsmath}
\usepackage{esint}
\usepackage{amssymb}
\usepackage{textcomp}
\usepackage{gensymb}
\usepackage{array}
\usepackage{titlesec}
\usepackage[margin=1in]{geometry}

\setlength{\parindent}{0pt}

\newcommand{\msp}{\text{ }}
\newcommand{\ligneh}{\msp | \msp}
\newcommand{\si}{\text{ si }}
\newcommand{\et}{\text{ et }}
\newcommand{\ou}{\text{ ou }}
\newcommand{\ch}{\text{ch}}
\newcommand{\sh}{\text{sh}}
\newcommand{\rg}{\text{rg}}
\newcommand{\Vect}{\text{Vect}}
\newcommand{\dl}[2]{\text{DL}_{#1}(#2)}
\newcommand{\mat}{\text{Mat}}
\newcommand{\ind}{\perp\!\!\!\perp} 
\newcommand{\conv}{\text{ converge }}
\newcommand{\cov}{\text{Cov}}
\newcommand{\tr}{\text{tr}}
\newcommand{\nint}[2]{\llbracket {#1}, {#2} \rrbracket}
\newcommand{\cvu}{\overset{\text{CVU}}{\longrightarrow}}
\newcommand{\cvs}{\overset{\text{CVS}}{\longrightarrow}}
\newcommand{\cvn}{\overset{\text{CVN}}{\longrightarrow}}

\newcommand{\dd}{\text{d}}
\newcommand{\dx}{\text{d}x}
\newcommand{\dy}{\text{d}y}
\newcommand{\dz}{\text{d}z}
\newcommand{\dr}{\text{d}r}
\newcommand{\dt}{\text{d}t}
\newcommand{\dtheta}{\text{d}\theta}
\newcommand{\dphi}{\text{d}\phi}
\newcommand{\Dd}{\text{D}}
\newcommand{\repr}{|\mathcal{R}}
\newcommand{\grad}{\overrightarrow{{\text{grad}}} \;}
\newcommand{\bgrad}[1]{(\vec{#1}\cdot\overrightarrow{{\text{grad}}})}
\newcommand{\cste}{\text{cste}}
\newcommand{\rot}{\overrightarrow{\text{rot}} \;}

\setlength{\emergencystretch}{3em}
\hbadness=10000

\title{Formules}
\author{Huang Ziwen}
\date{2026-2027 PC*}

\makeatletter
\def\ps@myheader{
    \def\@evenhead{\hfil Formules de physique - Huang Ziwen - 2026-2027 PC* \hfil}%
    \def\@oddhead{\hfil Formules de physique - Huang Ziwen - 2026-2027 PC* \hfil}%
}
\makeatother

\pagestyle{myheader}

\begin{document}

\section{Incertitudes}

\begin{align*}
    s_{\text{exp}} &= \sigma_{N-1} = \sqrt{\frac{1}{N-1}\sum_{i=1}^{N}(m_i-\overline{m})^2} \\
    s_A &= \sqrt{\frac{1}{N}}s_{\text{exp}} \\
    s_B &= \frac{a}{\sqrt{3}}
\end{align*}

\subsection*{Propagation des incertitudes}

\[ s_B = \sqrt{\sum_{k=1}^{n}(\frac{\partial f}{\partial x_k}s(x_k))^2} \]

\section{Optique}

\subsection*{Formule de Cauchy}

\[ n = a + \frac{b}{\lambda^2} \]

\subsection*{Théorèmes de conjugaison}

\[\overline{FA} \cdot \overline{F'A'} = f'f=-f^2\]
\[ \frac{1}{\overline{OA}}-\frac{1}{\overline{OA'}} = -\frac{1}{f'} \]

\section{Électrocinétique}

\subsection*{Filtres}

\scriptsize
\begin{center}
    \renewcommand{\arraystretch}{2}
\begin{tabular}{ |c|c|c|c|c|c|m{5cm}| } 
     \hline
      & $ \underline{N} $ & $ \underline{D} $ & $\underline{H}$ usuel & Type & $ \phi $ & Remarques  \\ 
     \hline \hline
     $R[\underline{C}]$ & $ 1 $ & $ 1 + j(\frac{\omega}{\omega_c}) $ & $ \frac{1}{1 + jx} $ & Passe-bas & $ 0 \rightarrow -\frac{\pi}{2} $ & Pente HF: $-20$ dB par décade [Intégrateur] \\
     \hline
     $[\underline{R}]C$ & $ j(\frac{\omega}{\omega_c}) $ & $ 1 + j(\frac{\omega}{\omega_c}) $ & $ \frac{jx}{1 + jx} $ & Passe-haut & $ \frac{\pi}{2} \rightarrow 0 $ & Pente BF: $+20$ dB par décade [Dérivateur] \\
     \hline
     $RL[\underline{C}]$ & $ 1 $ & $ 1 - (\frac{\omega}{\omega_0})^2 + j\frac{1}{Q}\frac{\omega}{\omega_0} $ & $ \frac{1}{1 - x^2 +  j\frac{x}{Q}} $ & Passe-bas & $ 0 \rightarrow -\pi $ & Pente HF: $-40$ dB par décade [Double intégrateur]. \\
     \hline
     $R[\underline{L}]C$ & $ -(\frac{\omega}{\omega_0}) ^2$ & $ 1 - (\frac{\omega}{\omega_0})^2 + j\frac{1}{Q}\frac{\omega}{\omega_0} $ & $ \frac{-x^2}{1 - x^2 +  j\frac{x}{Q}} $ & Passe-haut & $ \pi \rightarrow 0 $ & Pente BF: $+40$ dB par décade [Double dérivateur]. \\
     \hline
     $[\underline{R}]LC$ & $ j\frac{1}{Q}\frac{\omega}{\omega_0} $ & $ 1 - (\frac{\omega}{\omega_0})^2 + j\frac{1}{Q}\frac{\omega}{\omega_0} $ & $ \frac{1}{1+jQ(x-\frac{1}{x})} $ & Passe-bande & $ \frac{\pi}{2} \rightarrow -\frac{\pi}{2} $ & Pente BF: $+20$ dB par décade [Dérivateur], Pente HF: $-20$ dB par décade [Intégrateur], Croisements des asymptotes varie selon $Q$  \\
     \hline
\end{tabular}
\end{center}
\normalsize

\clearpage

\section{Mécanique}

\subsection*{Coordonnées cylindriques}

\[ \frac{\dd \vec{e_r}}{dt}_{\repr} = \dot{\theta} \vec{e_{\theta}} \]
\[ \frac{\dd \vec{e_{\theta}}}{dt}_{\repr} = -\dot{\theta} \vec{e_{r}} \]
\[ \vec{v}_{\repr} = \dot{r} \vec{e_r} + r \dot{\theta} \vec{e_{\theta}} + \dot{z} \vec{e_z} \]
\[ \vec{a}_{\repr} = (\ddot{r} - r\dot{\theta}^2) \vec{e_r} + (r \ddot{\theta} + 2 \dot{r} \dot{\theta}) \vec{e_{\theta}} + \ddot{z} \vec{e_z} \]
\[ \grad f = \frac{\partial f}{\partial r} \vec{e_r} + \frac{1}{r} \frac{\partial f}{\partial \theta} \vec{e_\theta} + \frac{\partial f}{\partial z} \vec{e_z} \]

\subsection*{Coordonnées sphériques}

\[ \frac{\dd \vec{e_r}}{dt}_{\repr} = \dot{\theta} \vec{e_{\theta}} + \dot{\phi}\sin \theta \vec{e_{\phi}} \]
\[ \vec{v}_{\repr} = \dot{r} \vec{e_r} + r \dot{\theta} \vec{e_{\theta}} + r \dot{\phi} \sin \theta \vec{e_{\phi}} \]

\subsection*{Changements de base}
\[ \vec{v}_{\repr} (M) = \vec{v}_{\repr '} (M) + \vec{v}_{\repr} (O') + \vec{\Omega}_{\mathcal{R}' / \mathcal{R}} \wedge \overline{O'M} \]
\[ \vec{v}_{\repr} (M) = \vec{v}_{\repr '} (M) + \vec{v_{e}} (M) \]
\[ \vec{a}_{\repr}(M) = \vec{a}_{\repr '}(M) + 2 \vec{\Omega}_{\mathcal{R}' / \mathcal{R}}\wedge \vec{v}_{\repr '} (M) + \vec{a_e}(M) \]

\subsection*{Moment cinétique}
\[ \overrightarrow{L_{O\repr}}(M) = \overrightarrow{OM} \wedge \vec{p}_{\repr}(M) \]
\[ \vec{M}_O = \overrightarrow{OM} \wedge \vec{F} \]

\subsection*{Intéractions newtonniennes}

\[ \vec{e} = -\frac{mC}{K}\vec{v_{\repr}} - \vec{e_{\theta}} \]
\[ r = \frac{p}{-\text{sgn}(K) + e\cos{\theta}} \]
\[ E_{m\repr} = \frac{K^2}{2mC^2}(e^2-1) \]

\section{Thermodynamique}

\subsection*{Identités}
\[ \dd U = T \dd S - P \dd V \]
\[ \dd H = T \dd S + V \dd P \]
\[ \dd S = \frac{1}{T}\dd U + \frac{P}{T} \dd V \]

\subsection*{Gaz parfait}

\[ P = \frac{1}{3}n_Vm \langle v^2 \rangle \]

\subsection*{Cycles de Carnot}

\begin{align*}
    \eta &\leq 1 - \frac{T_f}{T_c} \\
    e_{\text{frigo}} &\leq \frac{T_f}{T_c - T_f} \\
    e_{\text{pac}} &\leq \frac{T_c}{T_c - T_f}
\end{align*}

\section{Mécanique des fluides}

\subsection*{Pression atmosphèrique}
\[ P(z) = P_0 e^{-\frac{M_{\text{mol}}gz}{RT}} = P_0 e^{-\frac{m_{\text{gaz}}gz}{kT}} \]

\subsection*{Nombre de Reynolds}

\[ R_e = \frac{\mu v L}{\eta} \]

\subsection*{Trainée d'un solide}

\[ C_x = \frac{f_{\text{trainée}}}{\frac{1}{2}\rho S v^2} \]

\begin{itemize}
    \item Pour $R_e < 1$, $C_x = \frac{24}{R_e}$, \[\vec{f} = -6\pi \eta r \vec{v}\]
    \item Pour $2\times10^3 < R_e < 2\times10^5$, $C_x = C$ constant, \[\vec{f} = - \frac{C}{2}\rho \pi r^2 v \vec{v}\]
\end{itemize}

\subsection*{Équation de Navier-Stokes}
\[ \mu \frac{\Dd \vec{v}}{\Dd t} = \vec{f_V} - \grad P + \eta \Delta \vec{v} \]

\subsection*{Équation d'Euler}
\[ \mu \left( \frac{\partial \vec{v}}{\partial t} + \frac{1}{2} \grad v^2 + (\rot \vec{v}) \wedge \vec{v} \right) = \vec{f}_V - \grad P \]

\section{Induction}

\subsection*{Solénoïde infini}
\[ \overrightarrow{B_{\text{int}}} = \mu_0 n I \vec{e_z} \]
\[ L = \mu_0 \frac{N^2}{L}S \]
\[ M = \mu_0 \frac{N_1 N_2}{l} S_2 \]

\section{Électromagnétisme}

\subsection*{Équations de Maxwell}
\begin{align*}
    \text{div } \vec{E} &= \frac{\rho}{\epsilon_0} \\
    \text{div } \vec{B} &= 0 \\
    \rot \vec{E} &= -\frac{\partial \vec{B}}{\partial t} \\
    \rot \vec{B} &= \mu_0 \vec{j} + \mu_0 \epsilon_0 \frac{\partial \vec{E}}{\partial t}
\end{align*}

\subsection*{Équation locale de Poynting}

\[ \frac{\partial}{\partial t} \left( \frac{1}{2}\epsilon_0 E^2 + \frac{B^2}{2\mu_0} \right) + \text{div}\left( \frac{\vec{E} \wedge \vec{B}}{\mu_0} \right) = - \vec{j} \cdot \vec{E} \]

\subsection*{Dipôles électriques}
\[ V = \frac{p \cos \theta}{4 \pi \epsilon_0 r^2} = \frac{\vec{p}\cdot\vec{e_r}}{4\pi\epsilon_0 r^2} \]
\[ \vec{E} = \frac{p}{4 \pi\epsilon_0r^3}(2 \cos \theta \vec{e_r} + \sin \theta \vec{e_{\theta}}) = \frac{3(\vec{p} \cdot \vec{r})\vec{r} - r^2\vec{p}}{4\pi\epsilon_0r^5} \]
\[ \vec{F} = \bgrad{p}_O \vec{E} \]
\[ \vec{M}_O = \vec{p} \wedge \vec{E}(O) \]
\[ E_p = - \vec{p} \cdot \vec{E}(O) \]

\section{Mécanique quantique}

\[ E = h\nu \]
\[ p = \frac{E}{c} = \frac{h}{\lambda} \]
\[ \lambda_{\text{DB}} = \frac{h}{p} \]

\subsection*{Principe d'indetermination}
\[ \Delta x \Delta p \geq \frac{\hbar}{2} \]

\section{Mathématiques}

\begin{itemize}
    \item Pythagore généralisé - $ a^2 = b^2 + c^2 - 2bc\cos{\theta} $
    \item Factorisation de cosinus - $ \cos{p} + \cos{q} = 2 \cos{\frac{p+q}{2}}\cos{\frac{p-q}{2}} $ 
    \item Dérivée partielle et taux d'accroissement - $f(x+\dx) - f(x) = \frac{\partial f}{\partial x} dx$
    \item Dérivée partielle - $\dd f(x, y) = \frac{\partial f}{\partial x} \dx + \frac{\partial f}{\partial y} \dy$
    \item Calcul de période d'un oscillation quelconque - $T = 2 \int_{-x_0}^{x_0} \frac{1}{\dot{x}} \dx = 4 \int_{0}^{x_0} \frac{1}{\dot{x}} \dx$
\end{itemize}

\subsection*{Équations différentielles}
\[ \dot{s}(t) + \frac{s(t)}{\tau} = f(t) \implies s(t) = Ae^{-\frac{t}{\tau}} + s_p(t) \] 
\[ \ddot{s}(t) + \omega_0^2s = \omega_0^2s_{eq} \implies s(t) = A\cos(\omega_0t)+B\sin(\omega_0t) + s_{eq} = D\cos(\omega_0t + \phi) + s_{eq} \]
\[ \ddot{s}(t) - \frac{1}{\tau ^ 2} s = \omega_0^2s_{eq} \implies s(t) = A e^{\frac{-t}{\tau}} + B e^{\frac{t}{\tau}} \]

\subsection*{Équation différentielle du 2nd ordre}
\[ \ddot{s}(t) + \frac{\omega_0}{Q}\dot{s}(t) + \omega_0^2s(t) = f(t) \]
\begin{itemize}
    \item $Q<\frac{1}{2}$ -
        \[ r = -\frac{\omega_0}{2Q} \pm \frac{\omega_0}{2Q}\sqrt{1-4Q^2} \]
        \[ s(t) = Ae^{r_1t} + Be^{r_2t} \]
    \item $Q=\frac{1}{2}$ -
        \[ r = -\frac{\omega_0}{2Q} = -\omega_0 \]
        \[ s(t) = (A + Bt)e^{\omega_0t} \]
    \item $Q>\frac{1}{2}$ -
        \[ r = -\frac{\omega_0}{2Q} \pm i\frac{\omega_0}{2Q}\sqrt{4Q^2-1} \]
        \[ \tau = \frac{2Q}{\omega_0}, \omega_p = \omega_0\sqrt{1-\frac{1}{4Q^2}} < \omega_0 \]
        \[ s(t) = (A\cos{(\omega_pt)} + B\sin{(\omega_pt)})e^{-\frac{t}{\tau}} \]
\end{itemize}

\subsection*{Relations particulières}
\[ \delta = |\ln{\frac{u(t) - u(t \rightarrow + \infty)}{u(t+T) - u(t \rightarrow + \infty)}}| = \frac{T}{\tau} = \frac{2\pi}{\sqrt{4Q^2-1}} \]
\[ Q = \frac{\omega_0}{\Delta \omega} \]

\end{document}